\chapter{PLANTEAMIENTO DEL PROBLEMA}
\label{ch:planteamiento}
\thispagestyle{empty}

Este capítulo delimita el problema práctico que originó la investigación y lo transforma en una cuestión técnica abordable mediante un artefacto computacional. Además, establece las preguntas, los objetivos, las proposiciones de diseño, las dimensiones operativas discretas y el procedimiento metodológico que conectan la heterogeneidad térmica del entorno con la verificación y evaluación de la solución en los capítulos siguientes.

\section{DIAGNÓSTICO}
\label{sec:diagnostico}

Las redes eléctricas deben atender el crecimiento sostenido de la demanda e integrar centrales de generación renovable distribuidas geográficamente dentro del sistema eléctrico. En América Latina y el Caribe, la demanda de generación eléctrica aumentó 5,5\,\pct{} en 2024 y las fuentes renovables aportaron 69\,\pct{} de ese crecimiento; dentro de esa participación, 45\,\pct{} correspondieron a centrales hidroeléctricas convencionales y 19\,\pct{} a generación eólica y solar \parencite{olade2025generacionALC2024}.

En el Perú, entre 2010 y 2025, la producción eléctrica nacional aumentó de \SI{35908}{\giga\watt\hour} a \SI{65109}{\giga\watt\hour}, lo que representó un crecimiento acumulado de 81\,\pct{} y una tasa media anual de 4\,\pct{}. En el mismo período, la máxima demanda del Sistema Eléctrico Interconectado Nacional (SEIN) se elevó de \SI{4579}{\mega\watt} a \SI{7942}{\mega\watt}, con un incremento acumulado de 73\,\pct{} y una tasa anualizada de 3,7\,\pct{}. Entre 2023 y 2025, la producción creció 4,9\,\pct{} y la máxima demanda lo hizo en 4,4\,\pct{} \parencites[sec.~10.4]{minem2025anuario2024}[cuadro~4]{minem2025indicadoresDiciembre}{coesMaximaDemanda}.

La longitud acumulada de las líneas troncales de 500 y \SI{220}{\kilo\volt} del SEIN pasó de \SI{8266}{\kilo\metre} en 2010 a \SI{17179}{\kilo\metre} en 2025, equivalente a un aumento de 108\,\pct{}. La Figura~\ref{fig:demanda-expansion-red-peru} compara este inventario con la evolución de la máxima demanda. Desde 2020, la longitud de la red troncal creció en menor proporción que la demanda máxima, lo cual muestra una mayor exigencia sobre la capacidad de transporte existente \parencites{coesEstadisticasAnuales}{coesMaximaDemanda}.

\begin{figure}[htbp]
\centering
\begin{tikzpicture}
\begin{axis}[
  name=expansion,
  width=0.88\textwidth,
  height=0.50\textwidth,
  xmin=2009.5, xmax=2025.5,
  ymin=7500, ymax=20000,
  xtick={2010,2011,2012,2013,2014,2015,2016,2017,2018,2019,2020,2021,2022,2023,2024,2025},
  scaled x ticks=false,
  xticklabel={\pgfmathprintnumber[1000 sep={}]{\tick}},
  xticklabel style={rotate=45,anchor=east,font=\footnotesize},
  ytick={7500,10000,12500,15000,17500,20000},
  scaled y ticks=false,
  yticklabel={\pgfmathprintnumber[1000 sep={\,}]{\tick}},
  yticklabel style={font=\footnotesize,text=unired},
  xlabel={Año},
  ylabel={Líneas de 500 y 220 kV (km)},
  ylabel style={text=unired,font=\footnotesize},
  xlabel style={font=\footnotesize},
  axis y line*=right,
  axis x line*=bottom,
  grid=major,
  grid style={gray!20},
  legend style={at={(0.02,0.87)},anchor=north west,font=\scriptsize,draw=none,fill=white,fill opacity=0.88,text opacity=1}
]
\pgfplotsextra{
  \foreach \x/\lo/\hi in {
    2010/7500/8266, 2011/8266/9751, 2012/9751/10383,
    2013/10383/11569, 2014/11569/12573, 2015/12573/13453,
    2016/13453/14139, 2017/14139/15509, 2018/15509/15782,
    2019/15782/16510, 2020/16510/16383, 2021/16383/16101,
    2022/16101/16350, 2023/16350/16781, 2024/16781/17133,
    2025/17133/17179}{
      \pgfmathsetmacro{\xl}{\x-0.27}
      \pgfmathsetmacro{\xr}{\x+0.27}
      \draw[fill=unired!35,draw=unired!75]
        (axis cs:\xl,\lo) rectangle (axis cs:\xr,\hi);
  }
}
\addlegendimage{area legend,fill=unired!35,draw=unired!75}
\addlegendentry{Variación anual de la longitud acumulada}
\end{axis}

\begin{axis}[
  at={(expansion.south west)},anchor=south west,
  width=0.88\textwidth,
  height=0.50\textwidth,
  xmin=2009.5, xmax=2025.5,
  ymin=4400, ymax=8200,
  ytick={4500,5000,5500,6000,6500,7000,7500,8000},
  yticklabel style={font=\footnotesize,text=darkblue},
  ylabel={Máxima demanda del SEIN (MW)},
  ylabel style={text=darkblue,font=\footnotesize},
  axis y line*=left,
  axis x line=none,
  xtick=\empty,
  legend style={at={(0.02,0.97)},anchor=north west,font=\scriptsize,draw=none,fill=white,fill opacity=0.88,text opacity=1},
  clip=false
]
\addplot[color=darkblue,mark=*,mark size=2.0pt,very thick] coordinates {
  (2010,4578.94) (2011,4961.19) (2012,5291.00) (2013,5575.24)
  (2014,5737.27) (2015,6275.00) (2016,6492.41) (2017,6559.06)
  (2018,6884.59) (2019,7017.57) (2020,7125.30) (2021,7173.03)
  (2022,7467.45) (2023,7605.51) (2024,7794.01) (2025,7942.00)
};
\addlegendentry{Máxima demanda}
\end{axis}
\end{tikzpicture}
\caption[Máxima demanda y longitud de la red troncal]{Evolución de la máxima demanda del SEIN y de la longitud acumulada de líneas troncales de 500 y \SI{220}{\kilo\volt}, 2010--2025. Fuente: Elaboración propia con datos de las estadísticas anuales y del reporte de máxima demanda del COES \parencites{coesEstadisticasAnuales}{coesMaximaDemanda}.}
\label{fig:demanda-expansion-red-peru}
\end{figure}

La presión sobre el sistema de transporte también se explica por el cambio en la matriz de generación. En 2025, las fuentes renovables aportaron cerca del 63\,\pct{} de la producción nacional: 52\,\pct{} correspondieron a centrales hidroeléctricas convencionales y 11\,\pct{} a recursos energéticos renovables (RER) no convencionales. Como muestra la Figura~\ref{fig:participacion-renovable-peru}, la participación conjunta solar y eólica creció desde un valor cercano a cero en 2010 hasta 9,5\,\pct{} en 2025, con un aumento de 4,2 puntos porcentuales entre 2023 y 2025 \parencites[sec.~10.4]{minem2025anuario2024}[cuadro~4]{minem2025indicadoresDiciembre}. Esta expansión requiere redes colectoras, conexiones y capacidad de transporte entre nuevos puntos del SEIN \parencites{coes2023diagnostico2025_2034}{iea2025transmissiongrid}.

\begin{figure}[htbp]
\centering
\begin{tikzpicture}
\begin{axis}[
  name=variacion,
  width=0.88\textwidth,
  height=0.50\textwidth,
  ylabel={Variación interanual (\pct{})},
  xmin=2009.5, xmax=2025.5,
  ymin=-0.5, ymax=3.1,
  ytick={-0.5,0,0.5,1,1.5,2,2.5,3},
  axis y line*=right,
  axis x line=none,
  xtick=\empty,
  ybar,
  bar width=7pt,
  tick label style={font=\footnotesize},
  label style={font=\footnotesize}
]
\addplot[fill=lightblue,draw=darkblue!60] coordinates {
  (2011,0.000)
  (2012,0.135)
  (2013,0.319)
  (2014,0.546)
  (2015,0.708)
  (2016,0.813)
  (2017,0.058)
  (2018,1.513)
  (2019,0.150)
  (2020,0.670)
  (2021,-0.343)
  (2022,0.038)
  (2023,0.726)
  (2024,2.788)
  (2025,1.386)
};
\end{axis}
\begin{axis}[
  at={(variacion.south west)},
  anchor=south west,
  width=0.88\textwidth,
  height=0.50\textwidth,
  ylabel={Participación en la generación nacional (\pct{})},
  xlabel={Año},
  xmin=2009.5, xmax=2025.5,
  xtick={2010,2012,2014,2016,2018,2020,2022,2024,2025},
  xticklabel={\pgfmathprintnumber[fixed,precision=0,1000 sep={}]{\tick}},
  xticklabel style={rotate=45,anchor=east,font=\footnotesize},
  ymin=0, ymax=11,
  ytick={0,2,4,6,8,10},
  axis y line*=left,
  axis x line*=bottom,
  grid=major,
  grid style={gray!25},
  axis line style={black!70,line width=0.8pt},
  tick label style={font=\footnotesize},
  label style={font=\footnotesize},
  legend style={font=\footnotesize,at={(0.02,0.98)},anchor=north west,draw=none,fill=white}
]
\addplot[color=darkblue,mark=*,mark size=2.2pt,very thick] coordinates {
  (2010,0.003)
  (2011,0.003)
  (2012,0.138)
  (2013,0.457)
  (2014,1.003)
  (2015,1.711)
  (2016,2.524)
  (2017,2.582)
  (2018,4.095)
  (2019,4.245)
  (2020,4.915)
  (2021,4.572)
  (2022,4.610)
  (2023,5.336)
  (2024,8.124)
  (2025,9.510)
};
\addlegendentry{Participación solar + eólica}
\addlegendimage{area legend,fill=lightblue,draw=darkblue!60}
\addlegendentry{Variación interanual}
\end{axis}
\end{tikzpicture}
\caption[Participación de la generación solar y eólica en el Perú]{Participación conjunta de la generación solar y eólica en la producción eléctrica nacional del Perú y variación interanual, 2010--2025. Notas: las barras expresan diferencias en puntos porcentuales, no tasas relativas; el valor de 2025 corresponde a cifras preliminares. Fuente: Elaboración propia con datos del Ministerio de Energía y Minas \parencites[sec.~10.4]{minem2025anuario2024}[cuadro~4]{minem2025indicadoresDiciembre}.}
\label{fig:participacion-renovable-peru}
\end{figure}

Los cables enterrados forman parte de las redes colectoras internas de las centrales eólicas y solares\footnote{XLPE corresponde a \textit{cross-linked polyethylene} o polietileno reticulado, un material aislante termoestable utilizado en cables de potencia \parencites[9--10]{thue1999electrical}{iec60502_2_2014}. La estructura multicapa de estos cables se describe en la Sección~\ref{subsec:generacion-evacuacion-calor} y se muestra en la Figura~\ref{fig:cable-multicapa}.}. Entre 2023 y 2025 se conectaron al SEIN siete centrales de al menos \SI{80}{\mega\watt}, cuatro eólicas y tres solares, que suman alrededor de \SI{1250}{\mega\watt} \parencites{osinergmin2024puntaLomitas}{minem2024nuevasRER}{minem2025nuevasCentrales}. La Central Eólica Punta Lomitas ilustra la magnitud de esta infraestructura: sus 57 aerogeneradores ---50 del proyecto original de \SI{260}{\mega\watt} y 7 de la expansión de \SI{36,4}{\mega\watt}--- se enlazan mediante circuitos colectores subterráneos de \SI{33}{\kilo\volt} \parencites{coes2021puntaLomitasEO}{engie2023puntaLomitas}. Los 15 circuitos iniciales totalizan unos \SI{75}{\kilo\metre} de trazado según el estudio de operatividad \parencite{coes2021puntaLomitasEO}, y la expansión añade dos circuitos adicionales \parencite{coes2022puntaLomitasExpansion}. Bajo una longitud media similar por circuito, los 17 circuitos equivalen a unos \SI{85}{\kilo\metre} de zanja y cerca de \SI{255}{\kilo\metre} de cable unipolar. El uso de cables enterrados también alcanza a enlaces urbanos de transmisión, como el circuito subterráneo de \SI{220}{\kilo\volt} previsto en la reconfiguración Chavarría--Santa Rosa--Carapongo en Lima \parencite[13]{coes2022planTransmision2023_2032} y los aproximadamente \SI{10}{\kilo\metre} de trazados subterráneos de \SI{60}{\kilo\volt} del proyecto Palián en Huancayo \parencites[168]{coes2022planTransmision2023_2032}[3,55--56]{proinversion2026palian}.

En estas instalaciones, la corriente máxima admisible o ampacidad ($\Imax$) no depende únicamente de la sección metálica del conductor. Esta capacidad resulta del equilibrio entre el calor generado y su disipación por las capas del cable, el lecho de asiento, el relleno y el suelo nativo \parencites{iec60287}{fariz2026}{mobius2025}. Los métodos clásicos y las normas proporcionan la base del cálculo, aunque suelen simplificar el entorno mediante resistencias térmicas equivalentes y propiedades homogéneas \parencites[752--760]{neher1957}{iec60287}{iec60853}. Esta simplificación puede perder precisión cuando la humedad, los estratos, el lecho de asiento, los rellenos térmicos o las zonas secas locales alteran la conductividad térmica alrededor del cable \parencites[4--6]{enescu2021}[1--3]{aldulaimi2024}.

La evidencia publicada confirma que las propiedades térmicas del entorno modifican de manera sustancial la capacidad de transporte calculada. La Tabla~\ref{tab:diagnostico-casos} y la Figura~\ref{fig:resistividad-ampacidad-casos} reúnen cuatro casos representativos en cables aislados con XLPE. En Sandton, Sudáfrica, elevar la resistividad de \SI{0,6}{\kelvin\metre\per\watt} a \SI{3,7}{\kelvin\metre\per\watt} redujo la ampacidad de \SI{518}{\ampere} a \SI{224}{\ampere} \parencite[5--6, tabs.~2--3]{khumalo2025}. En una planta fotovoltaica del sur del Perú, pasar de un suelo uniforme de \SI{1,5}{\kelvin\metre\per\watt} a las condiciones del sitio (\SI{4,9}{\kelvin\metre\per\watt} y temperatura local) redujo la ampacidad de \SI{563}{\ampere} a \SI{350}{\ampere} \parencite[7--8]{ariasVelasquez2023}. En un estudio peruano financiado por CONCYTEC para cables de \SI{220}{\kilo\volt}, variar la resistividad del entorno de \SI{2,5}{\kelvin\metre\per\watt} a \SI{3,0}{\kelvin\metre\per\watt} bajó la capacidad de \SI{1157}{\ampere} a \SI{1064}{\ampere} \parencite[15--16]{atoccsa2024}. Finalmente, en Corea del Sur, sustituir arena natural ($k=\SI{1,4}{\watt\per\metre\per\kelvin}$) por hormigón de agregado precolocado (PAC, $k=\SI{2,1}{\watt\per\metre\per\kelvin}$) elevó la ampacidad de aproximadamente \SI{941}{\ampere} a \SI{988}{\ampere} y redujo la temperatura máxima de \SI{78}{\celsius} a \SI{71}{\celsius} \parencite[12--16]{kim2025}.

\begin{table}[htbp]
\centering
\caption{Casos publicados que relacionan las propiedades térmicas del entorno con la ampacidad o la temperatura. Fuente: Elaboración propia a partir de las referencias indicadas.}
\label{tab:diagnostico-casos}
\begin{tabularx}{\textwidth}{L{2.7cm}L{3.2cm}Y Y}
\toprule
\textbf{Fuente} & \textbf{Configuración} & \textbf{Contraste estudiado} & \textbf{Resultado reportado} \\
\midrule
\textcite[5--6]{khumalo2025} & Cable tripolar XLPE de media tensión ($1\times3\times\SI{185}{\milli\metre\squared}$ Cu) & Resistividad de \SI{0,6}{\kelvin\metre\per\watt} frente a \SI{3,7}{\kelvin\metre\per\watt} en entierro directo & La ampacidad disminuyó de \SI{518}{\ampere} a \SI{224}{\ampere} ($-57\,\pct{}$). \\
\textcite[7--8]{ariasVelasquez2023} & Tres cables unipolares XLPE de \SI{33}{\kilo\volt} (\SI{630}{\milli\metre\squared} Al, trébol) & Suelo uniforme de \SI{1,5}{\kelvin\metre\per\watt} frente al valor reportado de \SI{4,9}{\kelvin\metre\per\watt} y temperatura de sitio & La ampacidad calculada disminuyó de \SI{563}{\ampere} a \SI{350}{\ampere} ($-38\,\pct{}$). \\
\textcite[15--16]{atoccsa2024} & Tres cables unipolares XLPE de \SI{220}{\kilo\volt} (\SI{2000}{\milli\metre\squared} Cu, plana) & Resistividad de \SI{2,5}{\kelvin\metre\per\watt} frente a \SI{3,0}{\kelvin\metre\per\watt} con relleno mejorado & El análisis de sensibilidad estimó una reducción de \SI{1157}{\ampere} a \SI{1064}{\ampere} ($-8\,\pct{}$). \\
\textcite[12--16]{kim2025} & Banco de seis cables XLPE de \SI{154}{\kilo\volt} (\SI{1200}{\milli\metre\squared} Cu, dos filas) & Hormigón de agregado precolocado (PAC, $k=\SI{2,1}{\watt\per\metre\per\kelvin}$) frente a arena natural ($k=\SI{1,4}{\watt\per\metre\per\kelvin}$) & La ampacidad aumentó aproximadamente de \SI{941}{\ampere} a \SI{988}{\ampere}; a igual corriente, $\Tmax$ bajó de \SI{78}{\celsius} a \SI{71}{\celsius}. \\
\bottomrule
\end{tabularx}
\end{table}

\begin{figure}[htbp]
\centering
\begin{tikzpicture}
\begin{groupplot}[
  group style={group size=4 by 1,horizontal sep=0.35cm},
  width=0.30\textwidth,
  height=0.52\textwidth,
  ymin=0, ymax=1600,
  xmin=0.05, xmax=0.95,
  xtick=\empty,
  ytick={0,200,400,600,800,1000,1200},
  grid=major,
  grid style={gray!25},
  axis line style={black!70},
  tick label style={font=\scriptsize},
  label style={font=\scriptsize},
  xlabel style={font=\tiny,align=center},
  clip=false
]
\nextgroupplot[
  ylabel={Ampacidad (\si{\ampere})},
  xlabel={\shortstack{Khumalo et al.\\MT; $1\times3\times\SI{185}{\milli\metre\squared}$\\Cu; tripolar\\entierro directo}}
]
\addplot[ybar,bar width=16pt,fill=darkblue!35,draw=darkblue] coordinates {(0.32,518)};
\addplot[ybar,bar width=16pt,fill=darkblue!85,draw=darkblue] coordinates {(0.68,224)};
\node[font=\scriptsize,align=center,anchor=north] at (axis cs:0.50,1580) {\textcolor{darkblue!55}{\rule{1.1ex}{1.1ex}} $\rhoth=0{,}6$\\$k=1{,}7$; 518 A};
\node[font=\scriptsize,align=center,anchor=north] at (axis cs:0.50,1350) {\textcolor{darkblue!90}{\rule{1.1ex}{1.1ex}} $\rhoth=3{,}7$\\$k=0{,}3$; 224 A\\$-57$\,\pct{} $I_{\max}$};

\nextgroupplot[
  yticklabels=\empty,
  xlabel={\shortstack{Arias Velásquez et al.\\\SI{33}{\kilo\volt}; $3\times\SI{630}{\milli\metre\squared}$\\Al; unipolares\\trébol}}
]
\addplot[ybar,bar width=16pt,fill=unired!30,draw=unired] coordinates {(0.32,563)};
\addplot[ybar,bar width=16pt,fill=unired!80,draw=unired] coordinates {(0.68,350)};
\node[font=\scriptsize,align=center,anchor=north] at (axis cs:0.50,1580) {\textcolor{unired!45}{\rule{1.1ex}{1.1ex}} $\rhoth=1{,}5$\\$k=0{,}7$; 563 A};
\node[font=\scriptsize,align=center,anchor=north] at (axis cs:0.50,1350) {\textcolor{unired!90}{\rule{1.1ex}{1.1ex}} $\rhoth=4{,}9$\\$k=0{,}2$; 350 A\\$-38$\,\pct{} $I_{\max}$};

\nextgroupplot[
  yticklabels=\empty,
  xlabel={\shortstack{Atoccsa et al.\\\SI{220}{\kilo\volt}; $3\times\SI{2000}{\milli\metre\squared}$\\Cu; unipolares\\plana}}
]
\addplot[ybar,bar width=16pt,fill=teal!30,draw=teal!80!black] coordinates {(0.32,1157)};
\addplot[ybar,bar width=16pt,fill=teal!80,draw=teal!80!black] coordinates {(0.68,1064)};
\node[font=\scriptsize,align=center,anchor=north] at (axis cs:0.50,1580) {\textcolor{teal!45}{\rule{1.1ex}{1.1ex}} $\rhoth=2{,}5$\\$k=0{,}4$; 1157 A};
\node[font=\scriptsize,align=center,anchor=north] at (axis cs:0.50,1420) {\textcolor{teal!90!black}{\rule{1.1ex}{1.1ex}} $\rhoth=3{,}0$\\$k=0{,}3$; 1064 A$^{*}$\\$-8$\,\pct{} $I_{\max}$};

\nextgroupplot[
  yticklabels=\empty,
  xlabel={\shortstack{Kim et al.\\\SI{154}{\kilo\volt}; $6\times\SI{1200}{\milli\metre\squared}$\\Cu; unipolares\\2 filas planas}}
]
\addplot[ybar,bar width=16pt,fill=violet!30,draw=violet!80!black] coordinates {(0.32,988)};
\addplot[ybar,bar width=16pt,fill=violet!80,draw=violet!80!black] coordinates {(0.68,941)};
\node[font=\scriptsize,align=center,anchor=north] at (axis cs:0.50,1580) {\textcolor{violet!45}{\rule{1.1ex}{1.1ex}} $\rhoth=0{,}48$\\$k=2{,}1$; 988 A};
\node[font=\scriptsize,align=center,anchor=north] at (axis cs:0.50,1350) {\textcolor{violet!90!black}{\rule{1.1ex}{1.1ex}} $\rhoth=0{,}73$\\$k=1{,}4$; 941 A$^{\dagger}$\\$-5$\,\pct{} $I_{\max}$};
\end{groupplot}
\end{tikzpicture}
\caption[Resistividad térmica y ampacidad en cuatro casos publicados]{Cambio de ampacidad entre dos valores de resistividad térmica en cuatro instalaciones diferentes. Nota: la conductividad térmica es la inversa de la resistividad, $k=1/\rhoth$; en las etiquetas, $\rhoth$ se expresa en \si{\kelvin\metre\per\watt} y $k$ en \si{\watt\per\metre\per\kelvin}. $^{*}$Aproximación local de sensibilidad. $^{\dagger}$Valor leído aproximadamente de la Figura~15 de la fuente, para septiembre y una temperatura del terreno de \SI{21}{\celsius}. Fuente: Elaboración propia con datos de \textcite[5--6, tabs.~2--3]{khumalo2025}, \textcite[7--8]{ariasVelasquez2023}, \textcite[15--16]{atoccsa2024} y \textcite[4,12--13, fig.~15]{kim2025}.}
\label{fig:resistividad-ampacidad-casos}
\end{figure}

La dificultad metodológica radica en seleccionar una representación adecuada del entorno térmico. Un modelo homogéneo sustituye la distribución espacial por una conductividad única equivalente $k_{eq}$, lo cual facilita el cálculo pero puede ocultar zonas locales de baja disipación. Los métodos de elementos finitos (FEM) conservan la geometría y los materiales con detalle espacial, aunque exigen construir, mallar y verificar cada escenario \parencites[3--4,93--99]{cigre2025}[1--3]{aldulaimi2024}.

La Figura~\ref{fig:brecha} ordena la cadena que conecta la presión de la demanda y de la generación renovable sobre el sistema eléctrico con el riesgo de decisiones térmicas imprecisas. Si el modelo subestima la temperatura del conductor, puede admitir corrientes excesivas y acelerar el envejecimiento del aislamiento \parencite[1--2,9]{khumalo2025}; si la sobreestima, puede restringir capacidad utilizable y promover el sobredimensionamiento \parencites[4--6]{enescu2021}[1--2,9]{khumalo2025}. Por ello, esta investigación separó el efecto físico de la heterogeneidad térmica del error numérico del artefacto mediante escenarios controlados y referencias convergentes.

\begin{figure}[htbp]
\centering
\resizebox{1.0\textwidth}{!}{%
\begin{tikzpicture}[
  scale=1.,
  transform shape,
  box/.style={draw,rounded corners=3pt,align=center,text width=3.25cm,minimum height=1.25cm,font=\tiny},
  arr/.style={-{Stealth[length=3mm]},thick},
  darr/.style={-{Stealth[length=3mm]},thick,dashed}]
\node[box,fill=lightgray] (cables) {Tramos con cables\\enterrados\\restricciones de espacio,\\servidumbre y seguridad};
\node[box,fill=lightblue,right=0.75cm of cables] (real) {Entorno térmico\\enterrado\\$\kx$ variable: humedad,\\estratos, secado, suelo de relleno};
\node[box,fill=red!10,right=0.75cm of real] (gap) {Efecto no\\cuantificado sobre\\$T(x,y)$, $\Delta \Tmax$ y $\Delta \Imax$};
\node[box,fill=yellow!15,below=1.0cm of cables] (context) {Presión de red\\SEIN, demanda, renovables\\continuidad y capacidad};
\node[box,fill=green!10,below=1.0cm of real] (simpl) {Representación\\homogénea\\$k=k_{eq}$ y resistencias\\térmicas equivalentes};
\node[box,fill=orange!12,below=1.0cm of gap] (impact) {Riesgo de decisión\\subestimar: envejecimiento\\sobreestimar: capacidad ociosa};
\draw[arr] (context) -- (cables);
\draw[arr] (cables) -- (real);
\draw[arr] (real) -- (gap);
\draw[arr] (real) -- (simpl);
\draw[arr] (simpl) -- (gap);
\draw[darr] (context.south) -- ($(context.south)+(0,-0.55cm)$)
  -- ($(impact.south)+(0,-0.55cm)$) -- (impact.south);
\draw[arr] (gap) -- (impact);
\draw[arr] (simpl) -- (impact);
\end{tikzpicture}
}%
\caption{Relación entre presión de red, uso de cables enterrados, entorno térmico heterogéneo, simplificación homogénea y riesgo de decisión. Fuente: Elaboración propia.}
\label{fig:brecha}
\end{figure}

El diagnóstico delimita así la brecha de ingeniería: en el dominio bidimensional estacionario, resulta necesario establecer de forma verificable cuándo y cuánto cambian $T(x,y)$, $\Tmax$ e $\Imax$ según el contraste, el patrón, la extensión y la proximidad de las heterogeneidades del terreno circundante. Esta brecha sustenta la identificación del problema y la evaluación del artefacto PINN.

\section{IDENTIFICACIÓN Y DESCRIPCIÓN DEL PROBLEMA DE ESTUDIO}
\label{sec:identificacion-problema}

\subsection{OBJETO DE ESTUDIO Y UNIDAD DE ANÁLISIS}

El objeto de estudio físico es el arreglo de cables--entorno térmico en instalaciones subterráneas. Está conformado por los cables multicapa, su disposición geométrica, el lecho de asiento, el suelo de relleno, el suelo nativo, las fuentes de calor y las condiciones de frontera \parencites[3--6]{enescu2021}{ieee442}[1,10]{kim2025}.

La unidad de análisis es una sección transversal bidimensional parametrizada de ese arreglo con cables de polietileno reticulado (XLPE). Esta representación permite estudiar disposiciones individuales, planas o en trébol y conservar la ubicación relativa de cada material. El interés no recae en un fabricante específico, sino en configuraciones documentadas y reproducibles.

El artefacto PINN constituye la unidad de intervención tecnológica y no se confunde con el objeto físico estudiado. Recibe la especificación del caso, aproxima el campo $T(x,y)$ en cada región y produce indicadores térmicos y operativos verificables frente a referencias independientes.

\subsection{PROBLEMA REAL, SUBYACENTE Y TÉCNICO}

El problema real se manifiesta cuando las decisiones de dimensionamiento o de operación de cables enterrados se apoyan en estimaciones térmicas que no representan adecuadamente la heterogeneidad del entorno. Si el cálculo subestima la temperatura, el conductor puede superar su límite térmico admisible; si la sobreestima, puede restringirse capacidad disponible de la red.

El problema subyacente corresponde a un vacío de conocimiento sobre el efecto cuantitativo de la heterogeneidad espacial. No se ha establecido de forma trazable cómo el contraste, el patrón, la extensión y la proximidad de $\kx$ alteran $T(x,y)$, $\Tmax$ y el punto caliente frente a un suelo homogéneo equivalente. Tampoco está delimitado bajo qué condiciones esa diferencia altera decisiones de ingeniería. Un promedio único de conductividad no responde esta cuestión porque la geometría, las interfaces y las discontinuidades del suelo redistribuyen el flujo de calor \parencites[4--6]{enescu2021}[3--4,93--99]{cigre2025}.

El problema técnico consiste en diseñar y evaluar una formulación computacional integrada que represente explícitamente $\kx$, aproxime el campo $T(x,y)$ y cuantifique su efecto sobre $\Tmax$ e $\Imax$. El artefacto debe documentar exactitud, consistencia física, robustez y reproducibilidad mediante referencias analíticas, soluciones manufacturadas, referencias FEM convergentes y configuraciones publicadas \parencites{iec60287}{iec60853}[3--4,93--99]{cigre2025}[686--690]{raissi2019}.

\subsection{JUSTIFICACIÓN}

\paragraph{Justificación práctica.}
La justificación práctica responde al problema real de reducir la incertidumbre térmica que origina decisiones inseguras o excesivamente conservadoras sobre la ampacidad. La identificación de configuraciones donde una propiedad homogénea oculta una restricción térmica local concuerda con las reducciones de capacidad por aumento de resistividad y secado reportadas por \textcite[1,6--9]{khumalo2025}, y apoya decisiones de diseño y verificación en circuitos con bajo margen térmico \parencite[14842,14851]{fariz2026}.

\paragraph{Justificación tecnológica.}
La justificación tecnológica corresponde a la construcción de un artefacto que integra ecuaciones diferenciales en derivadas parciales (PDE) y aprendizaje automático. Al penalizar los residuos físicos de la ecuación de calor y las condiciones de interfaz dentro de la función de pérdida, la PINN genera un campo térmico continuo y diferenciable sin depender exclusivamente de pares etiquetados generados por FEM \parencites[686--690]{raissi2019}[1--3,12]{aldulaimi2024}. La propuesta no pretende sustituir las normas IEC ni a FEM, sino delimitar bajo qué condiciones el enfoque alcanza exactitud, consistencia física y robustez verificables \parencite[93--99]{cigre2025}.

\paragraph{Justificación metodológica.}
El enfoque de ciencia del diseño (\textit{design science research}, DSR) vincula la necesidad práctica con la especificación, construcción, demostración y evaluación iterativa del artefacto \parencite[55--56]{peffers2007}. Además del prototipo situado, el estudio aporta un procedimiento reproducible de verificación y principios de diseño transferibles sobre interfaces multicapa y heterogeneidad térmica \parencite[342--350]{gregor2013}.

\paragraph{Justificación científica y académica.}
La investigación conecta el dimensionamiento térmico de cables enterrados con el aprendizaje científico informado por física. Aunque estudios previos demostraron la viabilidad de las PINN en problemas térmicos controlados \parencites[1--3,17]{xing2023}[1,3,16]{pan2025}, persiste la necesidad de evaluar su comportamiento en geometrías multicapa con alto contraste de conductividad, múltiples interfaces y fuentes internas concentradas \parencite[2--6]{ren2025}.

\subsection{ALCANCES Y LIMITACIONES}

\paragraph{Alcances.}
El alcance temático abarca la conducción de calor estacionaria en secciones transversales bidimensionales de cables XLPE enterrados, incorporando explícitamente el conductor, las capas concéntricas, el lecho de asiento, el suelo de relleno y el suelo nativo. El alcance metodológico comprende el diseño y la evaluación mediante DSR de un artefacto PINN contrastado frente a soluciones analíticas o manufacturadas, referencias FEM convergentes y geometrías de la literatura. La ejecución de este estudio cubre el período comprendido entre julio de 2026 y junio de 2027, con fecha de corte bibliográfico en junio de 2026, e incluye la arquitectura lógica, el prototipo implementado, el catálogo de escenarios, las pruebas automatizadas y los principios de diseño derivados.

\paragraph{Limitaciones.}
El modelo es estacionario y bidimensional, por lo que excluye el régimen transitorio y las variaciones axiales. Las pérdidas eléctricas se representan como fuente térmica por efecto Joule, $I^2R(T)$, con la resistencia del conductor referida a \SI{20}{\celsius} y corregida linealmente según la temperatura media del conductor \parencites[752--760]{neher1957}[cláusulas~4--5]{iec60287}. No se resuelven las ecuaciones de Maxwell completas: cada caso emplea resistencia en corriente continua (DC) o en corriente alterna (AC) efectiva documentada sin duplicar factores de efecto piel o proximidad, y excluye pérdidas dieléctricas, de pantallas y de armaduras.

Asimismo, el modelo no acopla envejecimiento dieléctrico, migración de humedad ni convección subterránea; en consecuencia, las propiedades térmicas del suelo y del relleno se fijan como atributos físicos de cada escenario. La evaluación se limita a verificación numérica y comparación bibliográfica controlada, sin validación experimental de campo ni instrumentación en instalaciones reales. Por ello, los resultados no se generalizan fuera de las geometrías, fuentes y contornos verificados, ni certifican corrientes de operación en campo. Tampoco se incluye evaluación económica de CAPEX, OPEX o retorno financiero; toda referencia de desempeño corresponde al ámbito computacional, reportando hardware, tolerancias, semillas e inicializaciones como metadatos secundarios de reproducibilidad \parencite[350--351]{gregor2013}.

\section{FORMULACIÓN DEL PROBLEMA}
\label{sec:formulacion-problema}

La matriz de consistencia del Anexo~\ref{anx:matriz-consistencia} articula los problemas con los objetivos, las proposiciones de diseño, las variables discretas y los criterios de evaluación. A continuación se formulan el problema general y los cuatro problemas específicos que guiaron el desarrollo y la verificación del artefacto.

\subsection{FORMULACIÓN DEL PROBLEMA GENERAL}

En el arreglo de cables--entorno térmico, la heterogeneidad espacial del suelo nativo, el lecho de asiento y el suelo de relleno puede modificar el campo de temperatura $T(x,y)$, la temperatura máxima $\Tmax$ y la ampacidad $\Imax$ respecto de una representación homogénea equivalente \parencites[4--6]{enescu2021}[3--4,93--99]{cigre2025}. Sin embargo, resulta complejo cuantificar esas diferencias de forma reproducible y establecer cuándo alteran las decisiones de diseño y operación. Por ello, se plantea la siguiente interrogante central: ¿cómo diseñar y evaluar un artefacto computacional basado en una red neuronal informada por física para este dominio? En concreto, se busca representar la conductividad espacialmente variable $\kx$, estimar $T(x,y)$, $\Tmax$ e $\Imax$, y cuantificar las diferencias frente a representaciones homogéneas equivalentes.

\subsection{FORMULACIÓN DE LOS PROBLEMAS ESPECÍFICOS}

Los problemas específicos descomponen la brecha general en cuatro necesidades secuenciales, articuladas según la lógica DSR ilustrada en la Figura~\ref{fig:articulacion-problemas-especificos}. La síntesis metodológica posterior y la formulación de principios de diseño pertenecen al cierre del ciclo DSR y no constituyen un problema específico adicional.

\begin{enumerate}[label=PE\arabic*.]
  \item La representación física requerida por el artefacto no estaba formalizada en términos de entradas, requisitos funcionales y criterios de calidad que ordenen el cable multicapa, las interfaces, las fuentes, los contornos y $\kx$ en escenarios heterogéneos. ¿Qué entradas físicas, requisitos funcionales y criterios de calidad debe incluir una especificación trazable del arreglo de cables--entorno térmico?
  \item No se disponía de un paquete reproducible de soluciones analíticas o manufacturadas, referencias FEM con convergencia y casos publicados que permitiera construir y verificar el artefacto dentro del dominio delimitado. ¿Cómo construir y verificar el artefacto mediante referencias independientes y clasificar su aceptación?
  \item No se había cuantificado de forma trazable cómo el contraste, el patrón, la extensión y la proximidad de las heterogeneidades modifican $\Tmax$, el campo $T(x,y)$ y la localización del punto caliente. ¿Cómo cambian dichas magnitudes térmicas cuando varían esas cuatro dimensiones de heterogeneidad?
  \item No se habían identificado las condiciones que producen la mayor discrepancia de ampacidad entre una representación 2D heterogénea verificada y una representación homogénea equivalente. ¿Qué configuraciones y parámetros generan las mayores diferencias de $\Imax$ entre ambas representaciones?
\end{enumerate}

\begin{figure}[htbp]
\centering
\resizebox{\textwidth}{!}{%
\begin{tikzpicture}[
  >=Latex,
  problema/.style={draw, rounded corners=3pt, minimum height=2.35cm, text width=4.15cm, align=center, font=\small, text=black, inner sep=6pt},
  insumo/.style={align=center, font=\small, text width=4.15cm},
  flujo/.style={-{Stealth[length=3mm]}, line width=1.1pt},
  dato/.style={font=\small\bfseries, fill=white, inner sep=1.5pt}
]
  \node[problema, fill=lightgray] (pe1) {\textbf{PE1}\\Representación física no formalizada};
  \node[problema, fill=lightblue, right=1.25cm of pe1] (pe2) {\textbf{PE2}\\Referencias de verificación no integradas};
  \node[problema, fill=green!10, right=1.25cm of pe2] (pe3) {\textbf{PE3}\\Evaluación de escenarios no trazable};
  \node[problema, fill=yellow!15, below=2.25cm of pe2] (pe4) {\textbf{PE4}\\Discrepancias de ampacidad no identificadas};

  \node[insumo, above=0.55cm of pe1] (i1) {Cable multicapa, interfaces, fuentes, contornos y $\kx$};
  \node[insumo, above=0.55cm of pe2] (i2) {Soluciones analíticas, manufacturadas, FEM y casos publicados};
  \node[insumo, above=0.55cm of pe3] (i3) {Contraste, patrón, extensión y proximidad};
  \node[insumo, below=0.55cm of pe4] (i4) {Representaciones homogénea y heterogénea 2D};

  \draw[-{Stealth[length=3mm]}, line width=1.1pt] (i1) -- (pe1);
  \draw[-{Stealth[length=3mm]}, line width=1.1pt] (i2) -- (pe2);
  \draw[-{Stealth[length=3mm]}, line width=1.1pt] (i3) -- (pe3);
  \draw[-{Stealth[length=3mm]}, line width=1.1pt] (i4) -- (pe4);

  \draw[flujo] (pe1) -- node[dato] {D1} (pe2);
  \draw[flujo] (pe2) -- node[dato] {D2} (pe3);
  \draw[flujo] (pe3.south) -- node[dato] {D3} (pe4.north);

  \node[problema, fill=red!10, right=1.25cm of pe4] (salida) {\textbf{Resultado lógico}\\Base para formular los objetivos específicos};
  \draw[flujo] (pe4) -- (salida);
\end{tikzpicture}%
}
\caption[Articulación de los problemas específicos]{Articulación secuencial de los cuatro problemas específicos de la investigación. Cada problema recibe los elementos indicados y produce una definición o evidencia necesaria para abordar el siguiente. La síntesis metodológica posterior pertenece al ciclo DSR y no se formula como un problema específico adicional. Fuente: Elaboración propia.}
\label{fig:articulacion-problemas-especificos}
\end{figure}

\section{OBJETIVOS}
\label{sec:objetivos}

Los objetivos traducen los problemas formulados en metas verificables. Cada objetivo genera un producto discreto que sirve de insumo al siguiente y permite verificar su grado de cumplimiento al término del Capítulo~\ref{ch:propuesta-desarrollo}.

\subsection{OBJETIVO GENERAL}

Diseñar y evaluar un artefacto computacional basado en una red neuronal informada por física (PINN) que represente el arreglo de cables--entorno térmico, estime $T(x,y)$, $\Tmax$ e $\Imax$, cuantifique el efecto de la heterogeneidad térmica espacial y documente su exactitud, su consistencia física y su robustez. La reproducibilidad se sustenta en un repositorio con el código fuente, los datos empleados y sus fuentes, de modo que cada resultado pueda reconstruirse.

\subsection{OBJETIVOS ESPECÍFICOS}

La Figura~\ref{fig:articulacion-objetivos-especificos} muestra la articulación secuencial de los cuatro objetivos específicos. En esta cadena, cada producto intermedio funciona simultáneamente como variable dependiente de su etapa y como componente heredado de la entrada en la etapa siguiente.

\begin{enumerate}[label=OE\arabic*.]
  \item Elaborar una especificación trazable de las entradas físicas, los requisitos funcionales y los criterios de calidad del artefacto para representar el arreglo de cables--entorno térmico, incluidos los dominios multicapa, las fuentes térmicas, las condiciones de contorno, las interfaces y la conductividad térmica espacialmente variable $\kx$. \textit{Indicador:} cobertura superior al 90\,\pct{} de los requisitos críticos identificados y trazabilidad requisito--prueba (Anexo~\ref{anx:matriz-consistencia}).
  \item Construir y verificar el artefacto a partir de la especificación trazable y de un conjunto de soluciones analíticas o manufacturadas, referencias FEM convergentes y casos publicados, clasificando su estado de aceptación según exactitud, consistencia física y robustez entre ejecuciones. \textit{Indicador:} error absoluto de temperatura máxima $E_T < \SI{0,1}{\kelvin}$ y $\mathrm{NRMSE} < 1\,\pct{}$ en la puerta estricta de verificación F2, junto con los umbrales de la campaña explícita ejecutada ($e_{\Tmax} \le 5\,\pct{}$, $\mathrm{NRMSE} \le 5\,\pct{}$, desequilibrio $\le 2\,\pct{}$, salto de $T \le \SI{0,1}{\kelvin}$ y salto de flujo $\le 1\,\pct{}$).
  \item Aplicar el artefacto verificado a un catálogo de escenarios controlados de contraste, patrón, extensión y proximidad de heterogeneidades del suelo nativo y de los suelos de relleno, y consolidar un expediente de evidencia sobre $\Tmax$, el campo $T(x,y)$ y la localización del punto caliente. \textit{Indicadores:} exactitud frente a FEM ($E_T < \SI{0,5}{\kelvin}$ y $\mathrm{NRMSE} < 2,5\,\pct{}$ por semilla), efecto físico pareado $\Delta T_{\max}$ y posición del punto caliente respecto del escenario homogéneo, con efecto mínimo de interés de \SI{1}{\kelvin} y la regla de incertidumbre de la Sección~\ref{subsec:relevancia-decision}.
  \item Utilizar el expediente de evidencia térmica y un catálogo de representaciones pareadas para elaborar un expediente comparativo de ampacidad entre la representación 2D heterogénea y las representaciones homogéneas equivalentes, e identificar las condiciones que generan mayor discrepancia. \textit{Indicadores:} ampacidad $\Imax$ con error relativo frente a FEM menor que 1\,\pct{} por semilla, cambio porcentual $\Delta I$, efecto mínimo de interés de 2\,\pct{} y margen térmico; juicio de expertos complementario en escala de 1 a 10 (Anexos~\ref{anx:matriz-consistencia} y~\ref{anx:protocolo-expertos}).
\end{enumerate}

\begin{figure}[htbp]
\centering
\resizebox{\textwidth}{!}{%
\begin{tikzpicture}[
  >=Latex,
  objetivo/.style={draw, rounded corners=3pt, minimum height=2.35cm, text width=4.15cm, align=center, font=\small, text=black, inner sep=6pt},
  insumo/.style={align=center, font=\small, text width=4.15cm},
  flujo/.style={-{Stealth[length=3mm]}, line width=1.1pt},
  dato/.style={font=\small\bfseries, fill=white, inner sep=1.5pt}
]
  \node[objetivo, fill=lightgray] (oe1) {\textbf{OE1}\\Especificar entradas, requisitos y criterios};
  \node[objetivo, fill=lightblue, right=1.25cm of oe1] (oe2) {\textbf{OE2}\\Construir y verificar el artefacto};
  \node[objetivo, fill=green!10, right=1.25cm of oe2] (oe3) {\textbf{OE3}\\Evaluar escenarios heterogéneos};
  \node[objetivo, fill=yellow!15, below=2.25cm of oe2] (oe4) {\textbf{OE4}\\Comparar la ampacidad};

  \node[insumo, above=0.55cm of oe1] (i1) {\textbf{VI1: entrada inicial}\\paquete básico, completo o trazable};
  \node[insumo, above=0.55cm of oe2] (i2) {\textbf{Entrada nueva de VI2}\\referencia analítica/manufacturada, FEM, publicada o mixta};
  \node[insumo, above=0.55cm of oe3] (i3) {\textbf{Entrada nueva de VI3}\\catálogo base, estratificado, localizado, mejorado, seco o multimaterial};
  \node[insumo, below=0.55cm of oe4] (i4) {\textbf{Entrada nueva de VI4}\\par homogéneo equivalente, conservador u optimista};

  \draw[-{Stealth[length=3mm]}, line width=1.1pt] (i1) -- (oe1);
  \draw[-{Stealth[length=3mm]}, line width=1.1pt] (i2) -- (oe2);
  \draw[-{Stealth[length=3mm]}, line width=1.1pt] (i3) -- (oe3);
  \draw[-{Stealth[length=3mm]}, line width=1.1pt] (i4) -- (oe4);
  \draw[flujo] (oe1) -- node[dato] {VD1 $\equiv\mathrm{VI}_{2,h}$} (oe2);
  \draw[flujo] (oe2) -- node[dato] {VD2 $\equiv\mathrm{VI}_{3,h}$} (oe3);
  \draw[flujo] (oe3.south) -- node[dato] {VD3 $\equiv\mathrm{VI}_{4,h}$} (oe4.north);

  \node[objetivo, fill=red!10, right=1.25cm of oe4] (salida) {\textbf{Resultado}\\Expediente comparativo y evidencia para la síntesis DSR};
  \draw[flujo] (oe4) -- node[dato] {VD4} (salida);
\end{tikzpicture}%
}
\caption[Articulación de los objetivos específicos]{Articulación secuencial de los cuatro objetivos específicos. VI1 es la entrada inicial. En las etapas siguientes, cada VI es un paquete compuesto por la VD anterior ---entrada heredada--- y el objeto nuevo mostrado verticalmente: $\mathrm{VI2}=\{\mathrm{VD1},R_2\}$, $\mathrm{VI3}=\{\mathrm{VD2},E_3\}$ y $\mathrm{VI4}=\{\mathrm{VD3},C_4\}$. VD4 completa la evidencia sintetizada dentro de DSR. Fuente: Elaboración propia.}
\label{fig:articulacion-objetivos-especificos}
\end{figure}

\subsection{HIPÓTESIS Y PROPOSICIONES DE DISEÑO}
\label{subsec:proposiciones-diseno}

Conforme a la Guía N.° 02 de la Escuela de Posgrado de la Universidad Nacional de Ingeniería para maestrías de especialización o profesionalizantes, la tesis no realiza contrastación estadística inferencial de hipótesis sobre muestras poblacionales. No obstante, se conservan las proposiciones HE1--HE4 heredadas del plan de tesis como criterios de aceptación y proposiciones de diseño contrastables mediante evidencia computacional, sin realizar inferencias probabilísticas a partir de semillas aleatorias. La proposición general vincula la representación explícita de $\kx$ con la reproducción verificada de referencias y con la cuantificación de diferencias térmicas y de ampacidad frente a representaciones homogéneas equivalentes. La Figura~\ref{fig:articulacion-hipotesis-especificas} resume la articulación entre las cuatro proposiciones y sus indicadores.

\begin{enumerate}[label=HE\arabic*.]
  \item Un paquete físico inicial completo y documentado permite obtener VD1: especificación trazable aceptada que separe PDE, contornos, interfaces, fuentes y propiedades térmicas, cubra más del 90\,\pct{} de los requisitos críticos y vincule cada requisito con una prueba física o funcional.
  \item La especificación trazable y las referencias válidas permiten obtener VD2: artefacto aceptado dentro de su dominio de verificación. En la matriz del plan, la puerta F2 exige $E_T < \SI{0,1}{\kelvin}$, $\mathrm{NRMSE} < 1\,\pct{}$ y desequilibrio energético no mayor a 2\,\pct{}. En la campaña explícita del Capítulo~\ref{ch:propuesta-desarrollo}, la confirmación exige además $e_{\Tmax} \le 5\,\pct{}$, $\mathrm{NRMSE} \le 5\,\pct{}$, balances globales y regionales $\le 2\,\pct{}$, salto $\Delta T \le \SI{0,1}{\kelvin}$ y salto RMS de flujo $\le 1\,\pct{}$.
  \item El artefacto verificado permite obtener VD3: expediente térmico concluyente. A igualdad de geometría, pérdidas y contornos, se espera un aumento de $\Tmax$ ante baja conductividad cercana y una reducción ante relleno mejorado. Este efecto se evalúa en pares controlados mediante la regla $\delta_T = \SI{1}{\kelvin}$ y los criterios F3 ($E_T < \SI{0,5}{\kelvin}$ y $\mathrm{NRMSE} < 2,5\,\pct{}$).
  \item El expediente térmico y los pares de representación permiten obtener VD4: expediente comparativo de ampacidad concluyente. Empleando la media ponderada por área del entorno (Sección~\ref{subsec:homogenizacion}), se espera que promediar una zona local de baja conductividad próxima al conductor sobreestime $\Imax$. El contraste utiliza el umbral $\delta_I = 2\,\pct{}$, y la ausencia de reducción o una tendencia contraria refutarán esta expectativa.
\end{enumerate}

\begin{figure}[htbp]
\centering
\resizebox{\textwidth}{!}{%
\begin{tikzpicture}[
  >=Latex,
  hipotesis/.style={draw, rounded corners=3pt, minimum height=2.35cm, text width=4.15cm, align=center, font=\small, text=black, inner sep=6pt},
  indicador/.style={align=center, font=\small, text width=4.15cm},
  flujo/.style={-{Stealth[length=3mm]}, line width=1.1pt}
]
  \node[hipotesis, fill=lightgray] (he1) {\textbf{HE1}\\VD1: especificación trazable aceptada};
  \node[hipotesis, fill=lightblue, right=1.25cm of he1] (he2) {\textbf{HE2}\\VD2: artefacto aceptado en su dominio};
  \node[hipotesis, fill=green!10, right=1.25cm of he2] (he3) {\textbf{HE3}\\VD3: expediente térmico concluyente};
  \node[hipotesis, fill=yellow!15, below=2.25cm of he2] (he4) {\textbf{HE4}\\VD4: expediente comparativo clasificado};

  \node[indicador, above=0.55cm of he1] (m1) {Cobertura de requisitos y trazabilidad requisito--prueba};
  \node[indicador, above=0.55cm of he2] (m2) {$E_T$, $e_{\Tmax}$, NRMSE, desequilibrio energético y saltos};
  \node[indicador, above=0.55cm of he3] (m3) {$\Delta\Tmax$, punto caliente y campo $T(x,y)$};
  \node[indicador, below=0.55cm of he4] (m4) {$\Imax$, $\Delta I\pct$ y margen térmico};
  \draw[-{Stealth[length=3mm]}, line width=1.1pt] (m1) -- (he1);
  \draw[-{Stealth[length=3mm]}, line width=1.1pt] (m2) -- (he2);
  \draw[-{Stealth[length=3mm]}, line width=1.1pt] (m3) -- (he3);
  \draw[-{Stealth[length=3mm]}, line width=1.1pt] (m4) -- (he4);
  \draw[flujo] (he1) -- (he2);
  \draw[flujo] (he2) -- (he3);
  \draw[flujo] (he3.south) -- (he4.north);
  \node[hipotesis, fill=red!10, right=1.25cm of he4] (salida) {\textbf{Contraste global}\\Exactitud, efecto térmico y ampacidad};
  \draw[flujo] (he4) -- (salida);
\end{tikzpicture}%
}
\caption[Articulación de las proposiciones de diseño]{Articulación de las cuatro proposiciones de diseño y sus principales indicadores de contraste. La secuencia vincula la calidad del artefacto con la evidencia térmica y operativa; la reproducibilidad y la síntesis de principios se evalúan dentro del ciclo DSR. Fuente: Elaboración propia.}
\label{fig:articulacion-hipotesis-especificas}
\end{figure}

Los estados posibles de cada proposición comprenden aceptación, aceptación restringida, rechazo e inconclusión. La incertidumbre numérica y la variabilidad entre semillas deben ser inferiores al efecto interpretado, y no se considera cumplido OE4 mediante simulaciones térmicas a corriente o pérdida fija. Se distingue estrictamente $\Delta T_{\max}=T_{\max}^{\mathrm{het}}-T_{\max}^{\mathrm{hom}}$ (efecto físico entre escenarios) de $E_T=|T_{\max}^{\mathrm{PINN}}-T_{\max}^{\mathrm{ref}}|$ (error numérico frente a la referencia), en concordancia con los criterios de verificación de \textcite[93--99]{cigre2025} y la trazabilidad DSR de \textcite[350--351]{gregor2013}.

\subsection{DIMENSIONES OPERATIVAS DEL ESTUDIO}
\label{sec:dimensiones-operativas}

En coherencia con DSR, la operacionalización se estructura como una cadena discreta de construcción del artefacto. En cada etapa se distingue un paquete discreto de entrada (variable independiente, VI), la actividad de un objetivo específico y un producto verificable discreto (variable dependiente, VD). Salvo VD4, cada salida pasa a integrar la entrada siguiente: $\mathrm{VI1}=P_1$, $\mathrm{VI2}=\{\mathrm{VD1},R_2\}$, $\mathrm{VI3}=\{\mathrm{VD2},E_3\}$ y $\mathrm{VI4}=\{\mathrm{VD3},C_4\}$. Aquí $P_1$ denota el paquete de requisitos, $R_2$ las referencias, $E_3$ el catálogo térmico y $C_4$ los pares comparativos. En consecuencia, $\mathrm{VD1}\equiv\mathrm{VI}_{2,h}$, $\mathrm{VD2}\equiv\mathrm{VI}_{3,h}$ y $\mathrm{VD3}\equiv\mathrm{VI}_{4,h}$.

Esta definición separa niveles metodológicos. Tanto los atributos físicos ($\kx$, contraste $C_k$, extensión $f_h$, proximidad $d_h/D$, geometría, fuentes y contornos) como los hiperparámetros numéricos se fijan como controles o metadatos de ejecución, sin tratarse como variables estadísticas. La exactitud, la consistencia física, la robustez, la trazabilidad y la reproducibilidad actúan como criterios transversales DSR sobre VD1--VD4.

La Figura~\ref{fig:modelo-conceptual} representa el esquema entrada--proceso--salida de las dimensiones operativas y la Tabla~\ref{tab:dimensiones-operativas} especifica sus realizaciones, estados posibles e indicadores, manteniendo equivalencia directa con la matriz de consistencia (Anexo~\ref{anx:matriz-consistencia}) y la matriz de operacionalización (Anexo~\ref{anx:operacionalizacion}).

\begin{figure}[htbp]
\centering
\resizebox{0.94\textwidth}{!}{%
\begin{tikzpicture}[
  etapa/.style={draw,rounded corners=3pt,align=center,text width=3.15cm,minimum height=1.05cm,font=\small},
  io/.style={draw,rounded corners=3pt,align=center,text width=3.15cm,minimum height=0.95cm,font=\small},
  arr/.style={-{Stealth[length=3mm]},thick}]
\node[io,fill=lightgray] (vi1) {\textbf{VI1}\\Paquete inicial $P_1$};
\node[io,fill=lightblue,right=0.65cm of vi1] (vi2) {\textbf{VI2}\\$\{\mathrm{VD1},R_2\}$};
\node[io,fill=green!10,right=0.65cm of vi2] (vi3) {\textbf{VI3}\\$\{\mathrm{VD2},E_3\}$};
\node[io,fill=yellow!15,right=0.65cm of vi3] (vi4) {\textbf{VI4}\\$\{\mathrm{VD3},C_4\}$};
\node[etapa,below=0.55cm of vi1] (oe1) {\textbf{OE1}\\Especificar};
\node[etapa,below=0.55cm of vi2] (oe2) {\textbf{OE2}\\Construir y verificar};
\node[etapa,below=0.55cm of vi3] (oe3) {\textbf{OE3}\\Evaluar escenarios};
\node[etapa,below=0.55cm of vi4] (oe4) {\textbf{OE4}\\Comparar ampacidad};
\node[io,fill=lightgray,below=0.55cm of oe1] (vd1) {\textbf{VD1}\\Especificación trazable};
\node[io,fill=lightblue,below=0.55cm of oe2] (vd2) {\textbf{VD2}\\Artefacto verificado};
\node[io,fill=green!10,below=0.55cm of oe3] (vd3) {\textbf{VD3}\\Expediente térmico};
\node[io,fill=yellow!15,below=0.55cm of oe4] (vd4) {\textbf{VD4}\\Expediente de ampacidad};
\node[io,fill=red!8,text width=14.5cm,below=1.05cm of vd2,xshift=1.9cm] (p) {\textbf{Parámetros de caso y controles (no son VI):} $\kx$, $C_k$, $f_h$, $d_h/D$, geometría, pérdidas, contornos, dominio y configuración numérica aceptada.};
\foreach \n in {1,2,3,4}{\draw[arr] (vi\n) -- (oe\n); \draw[arr] (oe\n) -- (vd\n);}
\draw[arr] (vd1.east) to[bend left=22] node[font=\scriptsize,fill=white] {hereda} (vi2.south west);
\draw[arr] (vd2.east) to[bend left=22] node[font=\scriptsize,fill=white] {hereda} (vi3.south west);
\draw[arr] (vd3.east) to[bend left=22] node[font=\scriptsize,fill=white] {hereda} (vi4.south west);
\end{tikzpicture}
}%
\caption{Modelo entrada--proceso--salida de las dimensiones operativas. Las VI y las VD son objetos discretos; VD1--VD3 se heredan como componentes obligatorios de VI2--VI4. Los parámetros físicos y numéricos se documentan por caso, pero no se clasifican como VI ni VD. Fuente: Elaboración propia.}
\label{fig:modelo-conceptual}
\end{figure}

\begin{table}[htbp]
\centering
\caption{Variables discretas de entrada y salida, composición y estados por etapa de construcción del artefacto. Fuente: Elaboración propia.}
\label{tab:dimensiones-operativas}
\begingroup
\scriptsize
\setlength{\tabcolsep}{3pt}
\renewcommand{\arraystretch}{1.35}
\begin{tabularx}{\textwidth}{L{1.05cm}L{4.25cm}L{2.05cm}L{4.25cm}Y}
\toprule
\textbf{Etapa} & \textbf{VI discreta: composición y realizaciones} & \textbf{Actividad} & \textbf{VD discreta: objeto y estados} & \textbf{Indicadores del objeto} \\
\midrule
OE1 & \textbf{VI1=$P_1$. Paquete inicial de requisitos y datos.} Realizaciones: básico; completo; completo documentado. & Elaborar la especificación. & \textbf{VD1. Especificación trazable.} Estados: incompleta; completa pendiente de trazabilidad; trazable aceptada. & Cobertura de requisitos; trazabilidad requisito--prueba; consistencia de unidades y supuestos. \\
OE2 & \textbf{VI2=$\{\mathrm{VD1},R_2\}$. Expediente de construcción y verificación.} $R_2$: analítico/manufacturado; FEM convergente; publicado; mixto. & Construir y verificar. & \textbf{VD2. Artefacto verificado y expediente de aceptación.} Estados: rechazado; aceptado con restricciones; aceptado en el dominio. & $E_T$; $e_{\Tmax}$; NRMSE; residuos; balance energético; robustez; dominio aceptado. \\
OE3 & \textbf{VI3=$\{\mathrm{VD2},E_3\}$. Expediente de evaluación térmica.} $E_3$: base; estratificado; localizado; mejorado; seco; multimaterial. & Evaluar escenarios. & \textbf{VD3. Expediente de evidencia térmica.} Estados: inconcluso; sin efecto relevante; efecto relevante no crítico; condición crítica. & $T(x,y)$; $\Tmax$; $\Delta\Tmax$; gradiente; punto caliente; cobertura de escenarios. \\
OE4 & \textbf{VI4=$\{\mathrm{VD3},C_4\}$. Expediente de comparación.} $C_4$: par con homogénea equivalente; conservadora; optimista. & Comparar ampacidad. & \textbf{VD4. Expediente comparativo de ampacidad.} Estados: sin discrepancia relevante; reducción; aumento; inconcluso/fuera del dominio. & $\Imax$; $\Delta I\pct$; margen térmico; casos críticos; cobertura de pares. \\
\bottomrule
\end{tabularx}
\endgroup
\end{table}

\section{METODOLOGÍA}
\label{sec:metodologia}

\subsection{TIPO, NIVEL Y DISEÑO DE LA INVESTIGACIÓN}

La investigación adopta un enfoque cuantitativo porque calcula y compara temperaturas, errores espaciales, residuos físicos, balances energéticos, tiempos de cómputo y variaciones de ampacidad. Es aplicada y tecnológica porque construye y evalúa un artefacto computacional para resolver un problema concreto de ingeniería. Por su alcance explicativo y comparativo, analiza mediante pares controlados cómo el contraste, el patrón, la extensión y la proximidad de las heterogeneidades de $\kx$ modifican $T(x,y)$, $\Tmax$ e $\Imax$. El diseño corresponde a la ciencia del diseño (DSR) con verificación computacional frente a referencias analíticas, manufacturadas, FEM y publicadas, sin campaña experimental de campo.

\subsection{ADAPTACIÓN DE LA CIENCIA DEL DISEÑO Y SÍNTESIS DSR}

El estudio adaptó las seis actividades propuestas por \textcite[55--56]{peffers2007} bajo evaluación formativa continua, tal como recomiendan \textcite[198--203]{delport2024}, de manera que ningún incremento avanzara a la etapa siguiente sin superar las pruebas físicas y numéricas previas. El Procedimiento~\ref{proc:dsr} detalla la función de cada actividad.

\begin{procedure}[htbp]
\caption{Proceso DSR adaptado a la investigación. Fuente: Elaboración propia a partir de \textcite{peffers2007} y \textcite{gregor2013}.}
\label{proc:dsr}
\begin{enumerate}[label=\textbf{A\arabic*.}]
  \item \textbf{Identificación y motivación:} delimitar el arreglo de cables--entorno térmico, las entradas físicas relevantes y la brecha entre $k_{eq}$ y $\kx$.
  \item \textbf{Objetivos de solución:} transformar los requisitos físicos en productos verificables y umbrales cuantificables de exactitud, consistencia física, utilidad y reproducibilidad.
  \item \textbf{Diseño y desarrollo:} definir la arquitectura lógica, implementar los incrementos del prototipo y registrar decisiones, contratos y versiones.
  \item \textbf{Demostración:} ejecutar referencias independientes y escenarios heterogéneos controlados para generar campos térmicos e indicadores operativos.
  \item \textbf{Evaluación:} comparar los resultados con las referencias y los criterios preestablecidos, clasificando aceptación, restricción o rechazo.
  \item \textbf{Comunicación:} documentar el artefacto, los límites de aplicabilidad, el protocolo de reproducción y los principios de diseño derivados.
\end{enumerate}
\end{procedure}

Como cierre metodológico, la actividad de síntesis y comunicación DSR integra las evidencias producidas en VD1--VD4 para consolidar el protocolo reproducible y los principios de diseño delimitados por evidencia. Tal como ilustra la Figura~\ref{fig:sintesis-metodologica-dsr}, esta síntesis no añade un quinto objetivo ni una quinta hipótesis específica, sino que transforma los resultados evaluados en conocimiento comunicable \parencites[55--56]{peffers2007}[342--350]{gregor2013}.

\begin{figure}[htbp]
\centering
\resizebox{0.96\textwidth}{!}{%
\begin{tikzpicture}[
  evidencia/.style={draw,rounded corners=3pt,align=center,text width=3.35cm,minimum height=1.25cm,font=\small,inner sep=5pt},
  sintesis/.style={draw,rounded corners=3pt,align=center,text width=4.1cm,minimum height=1.55cm,font=\small,inner sep=6pt},
  salida/.style={draw,rounded corners=3pt,align=center,text width=3.7cm,minimum height=1.25cm,font=\small,inner sep=5pt},
  arr/.style={-{Stealth[length=2.8mm]},line width=1pt}
]
\node[evidencia,fill=lightgray] (e1) {VD1\\Especificación trazable};
\node[evidencia,fill=lightblue,below=0.45cm of e1] (e2) {VD2\\Artefacto verificado};
\node[evidencia,fill=green!10,below=0.45cm of e2] (e3) {VD3\\Expediente térmico};
\node[evidencia,fill=yellow!15,below=0.45cm of e3] (e4) {VD4\\Expediente de ampacidad};

\node[sintesis,fill=orange!12,right=1.5cm of e2,yshift=-0.85cm] (s) {\textbf{Actividad metodológica DSR}\\Evaluar, integrar, abstraer y comunicar la evidencia\\[1mm]\textit{No agrega problema, objetivo ni hipótesis}};

\node[salida,fill=red!8,right=1.5cm of s,yshift=1.05cm] (p1) {Procedimiento y guía de reproducción};
\node[salida,fill=red!8,right=1.5cm of s,yshift=-1.05cm] (p2) {Principios de diseño, límites y repositorio depurado};

\draw[arr] (e1.east) -- (s.west);
\draw[arr] (e2.east) -- (s.west);
\draw[arr] (e3.east) -- (s.west);
\draw[arr] (e4.east) -- (s.west);
\draw[arr] (s.east) -- (p1.west);
\draw[arr] (s.east) -- (p2.west);
\end{tikzpicture}%
}
\caption[Síntesis metodológica de la evidencia DSR]{Ubicación de la sistematización y de los principios de diseño dentro de la metodología DSR. La actividad ocurre después de obtener la evidencia de los cuatro objetivos específicos y no constituye un quinto problema, objetivo o hipótesis. Fuente: Elaboración propia.}
\label{fig:sintesis-metodologica-dsr}
\end{figure}

\subsection{ARTEFACTO PROPUESTO}

El artefacto es un modelo y prototipo computacional organizado en seis componentes lógicos, presentados en la Figura~\ref{fig:arquitectura}, con separación explícita entre la arquitectura lógica y su implementación. Las redes neuronales informadas por física se implementaron en PyTorch sobre Windows en precisión doble, mientras que las referencias numéricas de elementos finitos se generaron con FEniCSx y Gmsh en un entorno Linux bajo el subsistema WSL.

\begin{figure}[htbp]
\centering
\resizebox{0.96\textwidth}{!}{%
\begin{tikzpicture}[
  comp/.style={draw,rounded corners=3pt,align=center,text width=3.6cm,minimum height=1.25cm,font=\small},
  arr/.style={-{Stealth[length=2.8mm]},thick}]
\node[comp,fill=lightblue] (spec) {1. Especificación del caso\\geometría, materiales, carga y contornos};
\node[comp,fill=lightblue,right=1.0cm of spec] (phys) {2. Representación física\\PDE, fuentes, interfaces\\y conductividad variable $\kx$};
\node[comp,fill=green!12,right=1.0cm of phys] (solver) {3. Solucionador PINN\\aproximación y entrenamiento};
\node[comp,fill=orange!12,below=1.15cm of solver] (eval) {4. Verificación\\evaluación técnica con referencias, residuos, balance y robustez};
\node[comp,fill=orange!12,left=1.0cm of eval] (exp) {5. Experimentación numérica\\casos, sensibilidad y ampacidad};
\node[comp,fill=lightgray,left=1.0cm of exp] (report) {6. Trazabilidad\\datos, metadatos, resultados e informes};
\draw[arr] (spec) -- (phys);
\draw[arr] (phys) -- (solver);
\draw[arr] (solver) -- (eval);
\draw[arr] (eval) -- (exp);
\draw[arr] (exp) -- (report);
\draw[arr,dashed] (eval.north) to[bend left=25] (phys.south);
\draw[arr,dashed] (report.north) to[bend left=25] (spec.south);
\end{tikzpicture}%
}
\caption{Arquitectura lógica de seis componentes del artefacto DSR propuesto. Fuente: Elaboración propia.}
\label{fig:arquitectura}
\end{figure}

\subsection{POBLACIÓN, MUESTRA Y HOMOGENIZACIÓN DE PARES}
\label{subsec:homogenizacion}

La población de referencia corresponde al conjunto de configuraciones bidimensionales estacionarias de cables enterrados en entornos heterogéneos dentro del dominio delimitado. Al tratarse de una maestría de especialización, no se extrajo una muestra probabilística ni se realizó inferencia sobre poblaciones reales; en su lugar, se seleccionó intencionalmente un catálogo discreto de casos analíticos, manufacturados, publicados y escenarios controlados según su pertinencia técnica y capacidad de verificación.

Para contrastar el efecto de la heterogeneidad frente a una simplificación homogénea equivalente en el dominio externo a los cables $\Omega_e$ (suelo nativo, relleno e inclusiones), el comparador principal se define mediante el promedio de conductividad térmica ponderado por área de la Fórmula~\ref{for:keq-area}.

\begin{formula}[htbp]
\centering
\begin{equation}
\label{for:keq-area}
k_{eq}=\frac{1}{|\Omega_e|}\int_{\Omega_e}k(x,y)\,\mathrm{d}A
=\frac{\sum_m A_m k_m}{\sum_m A_m}.
\end{equation}
\caption{Conductividad térmica equivalente ponderada por área en el dominio externo. Fuente: Elaboración propia.}
\end{formula}

Las áreas $A_m$ de la Fórmula~\ref{for:keq-area} excluyen el interior de los cables y conservan idénticos la geometría, la profundidad, las capas internas, los contornos y la ley de pérdidas entre ambos miembros del par. El uso de la experimentación numérica en pares controlados dentro del ciclo DSR se organiza en la Tabla~\ref{tab:experimentacion-dsr}.

\begin{table}[htbp]
\centering
\caption{Uso de la experimentación numérica dentro del ciclo DSR. Fuente: Elaboración propia.}
\label{tab:experimentacion-dsr}
\begin{tabularx}{\textwidth}{L{3.2cm}Y Y}
\toprule
\textbf{Momento DSR} & \textbf{Propósito de la experimentación numérica} & \textbf{Evidencia esperada} \\
\midrule
Verificación formativa & Comprobar que las entradas físicas, PDE, contornos, interfaces, pérdidas y cálculo de ampacidad operan de forma coherente. & \textbf{VD2:} artefacto verificado, paquete de pruebas, errores frente a referencias, residuos, balance, dominio y estado de aceptación. \\
Demostración & Mostrar la transformación de escenarios con materiales heterogéneos en campos térmicos interpretables. & \textbf{VD3:} expediente térmico con mapas $T(x,y)$, $\Tmax$, punto caliente, diferencias, límites y estado del efecto. \\
Evaluación mediante verificación & Juzgar utilidad, robustez, reproducibilidad y límites frente a referencias independientes. & \textbf{VD4:} expediente comparativo con ampacidad, discrepancia, margen térmico, clasificación de casos y estado final. \\
\bottomrule
\end{tabularx}
\end{table}

\subsection{TÉCNICAS, INSTRUMENTOS Y MÉTRICAS DE EVALUACIÓN}
\label{subsec:relevancia-decision}

Los datos se recolectaron mediante revisión documental de fuentes primarias y registro sistemático de ejecuciones computacionales. Los instrumentos de control comprenden los contratos de caso en formato JSON, la matriz de extracción bibliográfica, la especificación de requisitos, las pruebas automatizadas unitarias y físicas, las mallas de referencia FEM, la bitácora DSR y los archivos de trazabilidad con huella criptográfica SHA-256 por corrida y caso.

Para evaluar cuantitativamente la exactitud, la consistencia física y el impacto operativo, la Fórmula~\ref{for:metricas-evaluacion} reúne las métricas operacionales adoptadas en el estudio, en concordancia con las recomendaciones de verificación de herramientas de cálculo térmico de \textcite[93--99]{cigre2025}.

\begin{formula}[htbp]
\centering
\begin{equation}
\label{for:metricas-evaluacion}
\begin{aligned}
E_T(\mathrm{K})&=|T_{\max}^{\mathrm{PINN}}-T_{\max}^{\mathrm{ref}}|,\\
\Delta T_{\max}(\mathrm{K})&=T_{\max}^{\mathrm{het}}-T_{\max}^{\mathrm{hom}},\\
e_{T_{\max}}(\pct)&=
100\frac{|T_{\max}^{\mathrm{PINN}}-T_{\max}^{\mathrm{ref}}|}
{T_{\max}^{\mathrm{ref}}-T_{\mathrm{amb}}},\\
\Delta I(\pct)&=100\frac{I_{\mathrm{het}}-I_{\mathrm{hom}}}{I_{\mathrm{hom}}},\\
\mathrm{NRMSE}(\pct)&=100\,\frac{\sqrt{\sum_{i=1}^{N}a_i\left(T_i^{\mathrm{PINN}}-T_i^{\mathrm{ref}}\right)^2}}{T_{\max}^{\mathrm{ref}}-T_{\mathrm{amb}}},\\
\varepsilon_E(\pct)&=100\,\frac{\left|P_{\mathrm{gen}}-P_{\mathrm{sal}}\right|}{P_{\mathrm{gen}}},\\
M_T(\mathrm{K})&=T_{\mathrm{lim}}-T_{\max}.
\end{aligned}
\end{equation}
\caption{Métricas operacionales de evaluación térmica, de ampacidad y de balance energético. Fuente: Elaboración propia.}
\end{formula}

En la Fórmula~\ref{for:metricas-evaluacion}, $E_T$ cuantifica el error absoluto de la temperatura máxima respecto a la referencia. Los indicadores $e_{T_{\max}}$ y $\mathrm{NRMSE}$ normalizan los errores mediante la elevación térmica sobre $T_{\mathrm{amb}}$, usando pesos $a_i>0$ ($\sum_i a_i=1$) para evitar dependencia del cero Celsius. El desequilibrio $\varepsilon_E$ contrasta la potencia generada $P_{\mathrm{gen}}=\int_\Omega Q\,\mathrm{d}A$ frente al flujo saliente $P_{\mathrm{sal}}=\int_{\partial\Omega}(-k\nabla T)\cdot\mathbf{n}\,\mathrm{d}s$. Por su parte, $\Delta T_{\max}$ y $\Delta I$ expresan los efectos físicos pareados, mientras $M_T$ registra el margen respecto al límite $T_{\mathrm{lim}}$. Adicionalmente, se reportan el error absoluto medio (MAE), el error máximo puntual y la exactitud relativa de ampacidad $e_I=100|I_{\max}^{\mathrm{PINN}}-I_{\max}^{\mathrm{FEM}}|/I_{\max}^{\mathrm{FEM}}$.

Como efectos mínimos de relevancia técnica se fijaron $\delta_T=\SI{1}{\kelvin}$ para $\Delta T_{\max}$ y $\delta_I=2\,\pct{}$ para $\Delta I$. Declarar un efecto pareado como resuelto exige que todas las semillas de confirmación de ambos miembros superen la aceptación y coincidan en signo con FEM. Además, tanto la dispersión entre semillas PINN como el cambio de malla FEM deben quedar por debajo de la magnitud calculada por FEM.

La Tabla~\ref{tab:criterios-evaluacion} consolida las reglas de aceptación del estudio. En ella se articulan las puertas de la matriz de consistencia (F1--F4), la puerta operativa de la campaña explícita A--D (Sección~\ref{sec:seleccion-pinn}) y el protocolo preliminar de juicio de expertos (Anexo~\ref{anx:protocolo-expertos}). La evaluación experta complementa la verificación numérica en cuanto a utilidad y claridad de límites, sin sustituir los criterios cuantitativos de error ni de conservación de la energía.

\begin{table}[htbp]
\centering
\caption{Criterios DSR de evaluación y reglas de aceptación del artefacto. Fuente: Elaboración propia a partir del plan de tesis.}
\label{tab:criterios-evaluacion}
\begin{tabularx}{\textwidth}{L{3.1cm}Y L{3.6cm}}
\toprule
\textbf{Criterio} & \textbf{Evidencia} & \textbf{Regla de aceptación} \\
\midrule
Especificación (F1) & Cobertura de requisitos críticos y matriz requisito--prueba. & Cobertura $>90\,\pct{}$ con trazabilidad verificada. \\
Verificación numérica (F2 y campaña A--D) & Concordancia de $\Tmax$ y del campo $T(x,y)$ frente a referencia convergente. & Matriz F2: $E_T<\SI{0,1}{\kelvin}$ y $\mathrm{NRMSE}<1\,\pct{}$. Campaña A--D: $e_{\Tmax}\le 5\,\pct{}$ y $\mathrm{NRMSE}\le 5\,\pct{}$. \\
Consistencia física & Residuos de PDE, contornos, saltos de interfaz y balance energético global y por región. & Desequilibrio $\varepsilon_E\le 2\,\pct{}$; salto térmico $\le\SI{0,1}{\kelvin}$; salto de flujo $\le 1\,\pct{}$. \\
Evaluación de escenarios (F3) & Concordancia en escenarios heterogéneos y resolución del efecto pareado frente a FEM. & $E_T<\SI{0,5}{\kelvin}$, $\mathrm{NRMSE}<2,5\,\pct{}$ y regla de relevancia ($\delta_T=\SI{1}{\kelvin}$). \\
Ampacidad (F4) & Búsquedas independientes por bisección acotada y contraste frente a FEM. & Error de $\Imax < 1\,\pct{}$ por semilla; relevancia $\delta_I=2\,\pct{}$. \\
Reproducibilidad & Reconstrucción desde código, datos, versiones, semillas y huellas SHA-256. & Repositorio trazable con todos los intentos registrados. \\
Juicio de expertos (F4) & Calidad técnica, consistencia, utilidad y límites evaluados frente a IEC y FEM. & Escala de 1 a 10 según protocolo preliminar (Anexo~\ref{anx:protocolo-expertos}). \\
\bottomrule
\end{tabularx}
\end{table}

\subsection{ETAPAS DE INTERVENCIÓN}

El Procedimiento~\ref{proc:etapas} organiza la intervención en siete fases operativas con su producto y su puerta de decisión, estableciendo la correspondencia con las fases F0--F4 de la matriz de consistencia.

\begin{procedure}[htbp]
\caption{Etapas operativas, productos y puertas de decisión de la investigación. Fuente: Elaboración propia a partir del plan de tesis.}
\label{proc:etapas}
\begin{tabularx}{\textwidth}{C{0.9cm}L{3.2cm}Y L{3.0cm}}
\toprule
\textbf{Fase} & \textbf{Actividad} & \textbf{Producto} & \textbf{Puerta} \\
\midrule
1 & Revisión y especificación (F0--F1) & VD1: especificación trazable y matriz requisito--prueba. & Aprobación de alcance y cobertura. \\
2 & Diseño lógico y prototipo (F2) & Arquitectura, modelo físico, código versionado y pruebas básicas. & Verificación unitaria y de operadores. \\
3 & Construcción, verificación y piloto (F2) & VD2: artefacto verificado mediante etapas A--D y expediente de aceptación. & Aceptación, restricción o rechazo. \\
4 & Demostración térmica (F3) & VD3: expediente de escenarios heterogéneos y clasificación de efectos. & Cobertura de pares y resolución. \\
5 & Comparación de ampacidad (F4) & VD4: expediente comparativo de corriente límite y casos críticos. & Validez operativa y juicio experto. \\
6 & Síntesis DSR & Protocolo reproducible, límites de validez y principios de diseño. & Evaluación final integrada. \\
7 & Comunicación & Tesis, repositorio depurado y documentación reproducible. & Revisión y defensa académica. \\
\bottomrule
\end{tabularx}
\end{procedure}

La ejecución de la fase de verificación comprendió etapas secuenciales de desarrollo y confirmación. Estas abarcaron selección de arquitectura (A), densidad de colocaciones (B), ajuste de aprendizaje (C), calibración de interfaz (C2) y confirmación con semillas independientes (D), detalladas en la Sección~\ref{sec:seleccion-pinn} del Capítulo~\ref{ch:propuesta-desarrollo}. Las referencias FEM conservaron exactamente la geometría, las conductividades, las fuentes y los contornos de cada problema PINN. Entre las dos últimas mallas se exigió variación de $\Tmax$ menor a \SI{0,1}{\kelvin} y 0,5\,\pct{}, cambio de RMSE regional menor a 0,5\,\pct{} y balance energético inferior a 0,2\,\pct{}.

El capítulo deja establecidos el problema, los objetivos, las proposiciones de diseño, la operacionalización discreta y el protocolo metodológico de evaluación. El Capítulo~\ref{ch:marco-teorico} reúne los antecedentes bibliográficos y las bases teóricas físicas, numéricas y metodológicas necesarias para fundamentar la construcción del artefacto y el análisis de sus resultados.
